\documentclass{amsart}
\usepackage{amsmath,amssymb,amsthm,hyperref}

\newtheorem{theorem}{Theorem}
\newtheorem{lemma}{Lemma}
\newtheorem{proposition}{Proposition}
\theoremstyle{definition}
\newtheorem{definition}{Definition}
\newtheorem{remark}{Remark}

\DeclareMathOperator{\Aut}{Aut}
\DeclareMathOperator{\Ham}{Ham}
\DeclareMathOperator{\Lie}{Lie}
\DeclareMathOperator{\grad}{grad}

\begin{document}

\title{No non-admissible extremal Sasaki metrics in the join
$M_3\star_{1,101}S^3_{3,2}$}
\maketitle

\section{Statement and answer}

Let $\Sigma_2$ be a closed Riemann surface of genus $2$, and let $M_3$ be the
total space of the primitive (Boothby--Wang) circle bundle over $\Sigma_2$
carrying the regular Sasaki structure whose transverse K\"ahler structure is
the hyperbolic metric of constant Gauss curvature $-1$ on $\Sigma_2$. Let
$S^3_{\mathbf w}$, $\mathbf w=(w_1,w_2)=(3,2)$, denote the $3$-sphere with the
weighted Sasaki structure whose transverse holomorphic structure is the
weighted projective line $\mathbb{CP}^1[3,2]$. Form the Sasaki join
\begin{equation}\label{eq:join}
  M \;=\; M_3 \star_{l_1,l_2} S^3_{\mathbf w},\qquad
  (l_1,l_2)=(1,101),
\end{equation}
a compact $5$-dimensional Sasaki manifold, with its distinguished join Sasaki
structure $(\mathcal S,\bar J)=(\xi,\eta,\Phi,g)$. Recall (see \S\ref{sec:bg})
that the \emph{admissible} Sasaki metrics in the Sasaki cone are those produced
by the join (admissible/Calabi-type) ansatz, i.e.\ those whose transverse
K\"ahler quotient is an \emph{admissible} K\"ahler metric on the
projectivization $P\to\Sigma_2$ compatible with the fibre--base splitting.

\begin{theorem}\label{thm:main}
On $M=M_3\star_{1,101}S^3_{3,2}$ \emph{every} extremal Sasaki metric lying in
the Sasaki class $(\mathcal S,\bar J)$ (equivalently, with fixed transverse
holomorphic structure of the join and Reeb field in the Sasaki cone
$\mathfrak t^+$) is admissible. Consequently the answer to the question is
\textbf{No}: there are no extremal Sasaki metrics in $(\mathcal S,\bar J)$ that
fail to be admissible.
\end{theorem}

The proof rests on two facts special to this example:
\begin{itemize}
\item[(A)] the genus of the base is $\ge 2$, so $\Sigma_2$ carries
\emph{no} non-zero holomorphic vector field; consequently \emph{every}
transversely holomorphic vector field on $M$ is vertical (tangent to the
$\mathbb{CP}^1[3,2]$ fibres) modulo the Reeb line;
\item[(B)] the maximal torus of the transverse automorphism group of
$(\mathcal S,\bar J)$ then has rank $2$ and the Sasaki cone is exactly the
$2$-dimensional join (``$\mathbf w$-Sasaki'') cone, every $T$-invariant
metric in which is admissible.
\end{itemize}
We recall the background and then prove Theorem~\ref{thm:main} via three lemmas.

\section{Background and conventions}\label{sec:bg}

\subsection{Sasaki structures and transverse K\"ahler geometry}
A Sasaki structure on a compact manifold $M^{2n+1}$ is a quadruple
$(\xi,\eta,\Phi,g)$ where $\xi$ is the Reeb (unit Killing) field, $\eta$ its
contact form, $\Phi$ an endomorphism with $\Phi^2=-\mathrm{Id}+\xi\otimes\eta$,
and $g$ the associated metric. The Reeb field generates a one-dimensional
foliation $\mathcal F_\xi$ whose transverse geometry is K\"ahler: on the normal
bundle $\nu(\mathcal F_\xi)=TM/\langle\xi\rangle$ one has a complex structure
$\bar J$ (induced by $\Phi$), a transverse K\"ahler form $\frac12 d\eta$, and a
transverse metric $g^T$. The \emph{transverse scalar curvature} $s^T$ is the
scalar curvature of $g^T$. A \emph{Sasaki class} is the set of Sasaki metrics
with fixed $(\xi,\bar J)$ and transverse K\"ahler class $[\frac12 d\eta]$,
deformed by transverse $\partial\bar\partial$-potentials.

\subsection{Extremal Sasaki metrics}
Following Boyer--Galicki--Simanca, a Sasaki metric is \emph{extremal} if the
transverse $(1,0)$-gradient $\partial^\#_g s^T$ of its transverse scalar
curvature is a transversely holomorphic vector field; equivalently $g$ is a
critical point of the transverse Calabi energy $\int_M (s^T)^2\,dV_g$ within
its Sasaki class. A standard identity (the transverse analogue of Calabi's)
shows this is equivalent to $\mathcal L_{\bar J\,\grad^T s^T}\,g^T=0$, i.e.\
$\bar J\,\grad^T s^T$ is a transverse Killing field.

\subsection{The Sasaki cone and the maximal torus}
Let $\Aut(\mathcal S,\bar J)$ be the automorphism group of the transverse
holomorphic structure (transverse biholomorphisms commuting with the flow of
$\xi$), with reduced/identity component $\Aut_0$ whose Lie algebra
$\mathfrak h$ consists of the transversely holomorphic vector fields commuting
with $\xi$. Let $T\subset\Aut_0$ be a maximal torus of rank $r$. The
(unreduced) \emph{Sasaki cone} is
\[
  \mathfrak t^+ \;=\; \{\,\xi'\in\Lie(T) : \eta(\xi')>0 \text{ everywhere}\,\},
\]
an open convex cone of dimension $r$; deforming $\xi$ within $\mathfrak t^+$
keeps $\bar J$ fixed (Boyer--Galicki--Simanca). The relevant search space for
extremal metrics in the problem is exactly $(\mathcal S,\bar J)$: Reeb field in
$\mathfrak t^+$ together with transverse K\"ahler potentials.

\subsection{The join and admissible metrics}
For \eqref{eq:join} the construction (Boyer--Tønnesen-Friedman) is the quotient
\[
  M_3\star_{l_1,l_2}S^3_{\mathbf w}
  =\bigl(M_3\times S^3_{\mathbf w}\bigr)/S^1_{(l_2,-l_1)},
\]
where $S^1_{(l_2,-l_1)}$ acts by the difference of the two Reeb circle actions
with the indicated weights. The quotient is Sasaki, with Reeb field a positive
combination of $\xi_{M_3}$ and $\xi_{\mathbf w}$. In the quasi-regular case the
transverse holomorphic quotient
\[
  P:=M/\mathcal F_\xi
\]
is the total space of a holomorphic $\mathbb{CP}^1[3,2]$ (weighted projective
line / log) bundle $\pi\colon P\to\Sigma_2$ with fibre the weighted projective
line $F:=\mathbb{CP}^1[3,2]$. The \emph{admissible} K\"ahler metrics on $P$ are
the Calabi-type metrics
\begin{equation}\label{eq:admissible}
  g^T \;=\; g_{\Sigma_2} \;\oplus\; \frac{d x^2}{\Theta(x)} \;\oplus\;
  \Theta(x)\,\theta\otimes\theta,\qquad x\in[-1,1],
\end{equation}
where $g_{\Sigma_2}$ is a fixed (hyperbolic) K\"ahler metric on $\Sigma_2$, $x$
the moment coordinate of the fibrewise $S^1$ action, $\theta$ a connection
$1$-form on the principal part, and $\Theta>0$ on $(-1,1)$ with
$\Theta(\pm1)=0$, $\Theta'(\pm1)=\mp2$, a smooth profile. Equivalently, the
admissible metrics are exactly the metrics that are invariant under the
fibrewise torus and split as base $\oplus$ fibre, with fibre data depending only
on the moment variable $x$.

\section{Three lemmas}

\begin{lemma}\label{lem:rigidbase}
$H^0(\Sigma_2,T^{1,0}\Sigma_2)=0$: $\Sigma_2$ admits no non-zero holomorphic
vector field, and $\Aut(\Sigma_2)$ is finite.
\end{lemma}

\begin{proof}
On a compact Riemann surface of genus $g$ the holomorphic vector fields are
$H^0(\Sigma,T^{1,0}\Sigma)=H^0(\Sigma,K_\Sigma^{-1})$, where $K_\Sigma$ is the
canonical bundle of degree $2g-2$, so $\deg K_\Sigma^{-1}=2-2g$. For $g=2$ this
is $2-4=-2<0$. A line bundle of negative degree on a compact Riemann surface
has no non-zero holomorphic section (a non-zero section would have an effective
divisor of degree equal to the bundle degree, which would be $\ge 0$). Hence
$H^0(\Sigma_2,T^{1,0}\Sigma_2)=0$. The identity component
$\Aut_0(\Sigma_2)$ is a complex Lie group with Lie algebra
$H^0(\Sigma_2,T^{1,0}\Sigma_2)=0$, hence trivial; together with Hurwitz's
finiteness of $\Aut(\Sigma_2)$ for $g\ge2$, $\Aut(\Sigma_2)$ is finite.
\end{proof}

\begin{lemma}\label{lem:vertical}
Every holomorphic vector field on the total space $P$ of the fibration
$\pi\colon P\to\Sigma_2$ (with rational fibre $F=\mathbb{CP}^1[3,2]$) is
vertical, i.e.\ tangent to the fibres:
$H^0(P,T_P)=H^0(P,T_{P/\Sigma_2})$. Consequently every transversely holomorphic
vector field on $M$ commuting with $\xi$ is, modulo the Reeb line, vertical.
\end{lemma}

\begin{proof}
Differentiating $\pi$ gives the short exact sequence of holomorphic sheaves
\[
  0 \longrightarrow T_{P/\Sigma_2} \longrightarrow T_P
  \xrightarrow{\ d\pi\ } \pi^*T_{\Sigma_2} \longrightarrow 0 ,
\]
where $T_{P/\Sigma_2}=\ker d\pi$ is the relative (vertical) tangent sheaf. Taking
global sections yields the left-exact sequence
\[
  0 \to H^0(P,T_{P/\Sigma_2}) \to H^0(P,T_P)
  \xrightarrow{\ d\pi\ } H^0(P,\pi^*T_{\Sigma_2}).
\]
By the projection formula,
$H^0(P,\pi^*T_{\Sigma_2})=H^0\!\bigl(\Sigma_2,T_{\Sigma_2}\otimes
\pi_*\mathcal O_P\bigr)$. Since the fibre $F$ is a compact connected rational
curve (weighted projective line), $H^0(F,\mathcal O_F)=\mathbb C$ and the higher
direct image relation $\pi_*\mathcal O_P=\mathcal O_{\Sigma_2}$ holds (the
fibres are connected with only constants as global functions). Hence
\[
  H^0(P,\pi^*T_{\Sigma_2})=H^0(\Sigma_2,T_{\Sigma_2})=0
\]
by Lemma~\ref{lem:rigidbase}. Therefore the map
$d\pi\colon H^0(P,T_P)\to H^0(P,\pi^*T_{\Sigma_2})$ is the zero map, and
$H^0(P,T_P)=H^0(P,T_{P/\Sigma_2})$: every holomorphic vector field on $P$ is
vertical.

For the Sasaki statement: a transversely holomorphic vector field on $M$
commuting with $\xi$ descends (in the quasi-regular case, and by continuity in
general) to a holomorphic vector field on $P=M/\mathcal F_\xi$, modulo the
Reeb generator $\xi$ itself. By the previous paragraph this descended field is
vertical, which is the assertion.
\end{proof}

\begin{lemma}\label{lem:cone}
For $M=M_3\star_{1,101}S^3_{3,2}$, the maximal torus of $\Aut_0(\mathcal S,\bar
J)$ has rank $2$, and the Sasaki cone $\mathfrak t^+$ is the $2$-dimensional
join cone, spanned by $\xi_{M_3}$ and $\xi_{\mathbf w}$. Moreover the
$T$-invariant transverse K\"ahler metrics with Reeb field in $\Lie(T)$ are
exactly the admissible metrics \eqref{eq:admissible}.
\end{lemma}

\begin{proof}
By Lemma~\ref{lem:vertical}, $\mathfrak h=\Lie\Aut_0(\mathcal S,\bar J)$ is
spanned by the Reeb direction $\xi$ together with the vertical (fibrewise)
holomorphic vector fields of $P\to\Sigma_2$. The fibre $F=\mathbb{CP}^1[3,2]$ is
a weighted projective line with distinct weights $w_1=3\neq w_2=2$; its
connected holomorphic automorphism group $\Aut_0(F)$ is the image of the
$\mathbb C^\times$ acting with weights $(3,2)$, i.e.\ a one-dimensional torus
$\mathbb C^\times$ fixing the two orbifold points $[1:0]$ and $[0:1]$ (the
``triangular'' unipotent flows present for the ordinary $\mathbb{CP}^1$ are
absent here because the two weights are distinct, so no extra $\mathfrak{sl}_2$
appears). Hence the vertical holomorphic automorphisms form a holomorphic
$\mathbb C^\times$-bundle over $\Sigma_2$ whose maximal compact subgroup is a
\emph{single} fibrewise circle $S^1_{\mathrm{fib}}$ of rank $1$; there is no
further compact toric direction because any additional commuting field would
have to be non-vertical, contradicting Lemma~\ref{lem:vertical}, or would
require an $\mathfrak{sl}_2$ in the fibre, which $\mathbb{CP}^1[3,2]$ does not
have.

Therefore a maximal torus $T$ of $\Aut_0(\mathcal S,\bar J)$ is generated by the
Reeb circle and $S^1_{\mathrm{fib}}$, of rank
\[
  r \;=\; 1\ (\text{Reeb}) \;+\; 1\ (\text{fibre}) \;=\; 2 .
\]
All maximal tori are conjugate, and the two generators are, up to the join
identification, $\xi_{M_3}$ and $\xi_{\mathbf w}$; hence
$\mathfrak t^+=\{\xi'\in\Lie(T):\eta(\xi')>0\}$ is the $2$-dimensional join
($\mathbf w$-Sasaki) cone.

Finally, suppose $g^T$ is a transverse K\"ahler metric invariant under $T$ with
Reeb field in $\Lie(T)$. The two commuting generators of $T$ are the Reeb field
and the fibre generator; the latter has a moment map which, in the fibre
$\mathbb{CP}^1[3,2]$, is a single coordinate $x\in[-1,1]$ (the image of the
fibrewise $S^1$ moment map). Invariance of $g^T$ under $S^1_{\mathrm{fib}}$
forces it to be constant along the $S^1_{\mathrm{fib}}$-orbits, so its fibre part
is determined by a function of $x$ alone; writing this function as a profile
$\Theta(x)$ and using the K\"ahler condition along the fibre gives precisely the
fibre block $\frac{dx^2}{\Theta(x)}\oplus\Theta(x)\theta\otimes\theta$ of
\eqref{eq:admissible}, with the smoothness/boundary conditions
$\Theta(\pm1)=0,\ \Theta'(\pm1)=\mp2$ imposed at the two fixed points
(orbifold poles). Along the base, $T$-invariance plus the absence of holomorphic
vector fields on $\Sigma_2$ (Lemma~\ref{lem:rigidbase}) leaves no freedom: the
horizontal block must be the pullback of the fixed base metric $g_{\Sigma_2}$,
up to a constant fixed by the cohomology class. Hence $g^T$ is admissible.
Conversely every admissible metric \eqref{eq:admissible} is manifestly
$T$-invariant with Reeb in $\Lie(T)$. This proves the stated equivalence.
\end{proof}

\begin{lemma}[Transverse Calabi symmetry]\label{lem:average}
Let $g$ be an extremal Sasaki metric in $(\mathcal S,\bar J)$ with Reeb field
$\xi\in\mathfrak t^+$. Then after a transverse biholomorphism the transverse
metric $g^T$ is invariant under the maximal torus $T$ of Lemma~\ref{lem:cone},
and the extremal field $\partial^\#_g s^T$ lies in
$\Lie(T)\otimes\mathbb C$.
\end{lemma}

\begin{proof}
This is the transverse Sasaki analogue of Calabi's theorem on the symmetry of
extremal K\"ahler metrics (Boyer--Galicki--Simanca; the proof transcribes
Calabi, ``Extremal K\"ahler metrics II,'' 1985, to the transverse K\"ahler
foliation). We verify its hypotheses and conclusion in the present setting.

\emph{Step 1: the extremal field generates transverse isometries.} By
extremality $V:=\partial^\#_g s^T$ is transversely holomorphic, and the
identity recalled in \S\ref{sec:bg} gives $\mathcal L_{\bar
J\grad^T s^T}g^T=0$, i.e.\ $X:=\bar J\,\grad^T s^T=\mathrm{Im}\,V$ is a
transverse Killing field. Its flow is a one-parameter group of transverse
isometries; its closure $T_V$ is a torus inside the isometry group
$\mathrm{Isom}(g^T)$ of the transverse metric, which is compact.

\emph{Step 2: maximal compact symmetry.} By the transverse Calabi theorem, the
identity component of $\mathrm{Isom}(g^T)$ for an extremal transverse K\"ahler
metric is a \emph{maximal compact subgroup} $K$ of $\Aut_0(\mathcal S,\bar J)$,
and the extremal field lies in the centre of $\Lie K$; in particular $V$ lies in
the complexified Lie algebra of a maximal torus $T'\subset K$. (The proof is the
foliated version of Calabi's: one shows that $\Aut_0$ is the complexification of
$\mathrm{Isom}_0(g^T)$ and that the extremal field is parallel to the centre.)

\emph{Step 3: identification of the torus.} By Lemma~\ref{lem:vertical} and the
computation in Lemma~\ref{lem:cone}, every maximal torus of
$\Aut_0(\mathcal S,\bar J)$ has rank $2$ (Reeb $+$ fibre), and all such tori are
conjugate. Conjugating by a transverse biholomorphism we arrange $T'=T$. Then
$g^T$ is $T$-invariant and $V=\partial^\#_g s^T\in\Lie(T)\otimes\mathbb C$.
\end{proof}

\section{Proof of Theorem~\ref{thm:main}}

Let $g$ be any extremal Sasaki metric in the class $(\mathcal S,\bar J)$, with
Reeb field $\xi\in\mathfrak t^+$. By Lemma~\ref{lem:average}, after a transverse
biholomorphism the transverse K\"ahler metric $g^T$ is invariant under the
maximal torus $T$ of Lemma~\ref{lem:cone}, and its Reeb field lies in $\Lie(T)$.

By the equivalence proved in Lemma~\ref{lem:cone}, the only $T$-invariant
transverse K\"ahler metrics with Reeb field in $\Lie(T)$ are the admissible
metrics \eqref{eq:admissible}: $T$-invariance forces the metric to be the fixed
base metric along $\Sigma_2$ (no freedom, by Lemma~\ref{lem:rigidbase}, since a
deformation along the base would be generated by a holomorphic vector field on
$\Sigma_2$, of which there are none), and to depend on the fibre only through the
single moment coordinate $x$ of the fibre torus, yielding the profile
$\Theta(x)$. Hence $g$ is admissible.

As $g$ was an arbitrary extremal metric in $(\mathcal S,\bar J)$, every extremal
Sasaki metric in this class is admissible. Equivalently, \emph{there are no}
extremal Sasaki metrics in $(\mathcal S,\bar J)$ that fail to be admissible.
This proves Theorem~\ref{thm:main}. \qed

\begin{remark}[On existence]
Theorem~\ref{thm:main} is a \emph{rigidity} statement: it says the extremal
locus contains no metric outside the admissible family; it does not assert that
extremal metrics exist. Within the admissible family the existence of extremal
representatives is governed by the positivity on $(-1,1)$ of an explicit
one-variable ``extremal polynomial'' (Boyer--Tønnesen-Friedman) whose
coefficients depend on $(l_1,l_2,\mathbf w)=(1,101,(3,2))$ and on the genus
constant of $\Sigma_2$. Whether or not that polynomial is positive, the
rigidity above holds.
\end{remark}

\begin{remark}[Role of the data]
The proof uses only two qualitative features of the data: the fibre is a
weighted projective line, giving a rank-$1$ fibre torus, and the base has genus
$\ge 2$, giving no holomorphic vector fields. Both hold here: $w_1=3\neq w_2=2$
are coprime, $101$ is prime and coprime to $w_1w_2=6$ (so the join is a smooth
quasi-regular Sasaki manifold of the stated type), and
$\mathrm{genus}(\Sigma_2)=2\ge 2$. The specific values $\mathbf w=(3,2)$ and
$l_2=101$ affect only the \emph{quantitative} admissibility analysis (topology
of the lens-space bundle $M$ and the coefficients of the extremal polynomial),
not the qualitative rigidity proved in Theorem~\ref{thm:main}.
\end{remark}

\end{document}
